\documentclass[10pt,a4paper]{amsart}
\usepackage[margin=19mm]{geometry}
\usepackage{amsmath,amssymb,mathtools}
\usepackage[scr=rsfs,cal=euler]{mathalfa}
\usepackage{xfrac}
\usepackage[unicode,hidelinks,bookmarks=false]{hyperref}
\newcommand{\ie}{i.e.\ }
\newcommand{\cf}{cf.\ }
\newcommand{\resp}{respectively\ }
\newcommand{\Subsection}{\subsection}
\providecommand{\nobreakdash}{\nobreak\hbox{-}}
\setlength{\emergencystretch}{2em}
\multlinegap0pt
\usepackage{mathtools}
\usepackage{amssymb}
\usepackage[shortlabels]{enumitem}
\newtheorem{theorem}{Theorem}
\newtheorem{lemma}{Lemma}
\newcommand{\E}{\mathbb{E}}
\newcommand{\cx}{\preceq_{\mathrm{cx}}}
\newcommand{\dd}{\,\mathrm{d}}
\newcommand{\eps}{\varepsilon}
\title{Convex order under majorization for projections of ell-p balls}
\author{Soufiane Fafe}
\date{Source: arXiv:2605.11451v2 (CC BY 4.0)}
\begin{document}
\maketitle

\noindent\emph{Source and licence.} Convex order under majorization for projections of ell-p balls. Version arXiv:2605.11451v2 (CC BY 4.0), distributed by arXiv under the Creative Commons Attribution 4.0 International licence, http://creativecommons.org/licenses/by/4.0/, as recorded in the arXiv licence metadata for this exact version and shown on the abstract page https://arxiv.org/abs/2605.11451v2. Full licence notice: \texttt{licenses/arxiv-2605.11451-CC-BY-4.0.txt}.\par\smallskip

\noindent\textit{Editorial scope.} Counted unit: the article's theorem thm:app-append (source0.tex lines 233--247). The article's complete original proof of it is delivered with it (source0.tex lines 1104--1216, one \texttt{proof} environment of 113 lines, which the article places away from the statement). No mathematics is written, repaired or re-derived: the statement and the proof are the article's own lines, in the article's own notation, and no retained line is substituted or re-flowed.  Retained uncounted as the source-local context the proof needs: lines 532--571; lines 550--558; lines 1061--1077.  Nothing was omitted from any retained span, so the package carries no omission record.  No retained line contains a citation, so the wrapper carries no bibliography and no bibliography entry.  No retained line contains a figure environment, an image-inclusion command or drawing code, so no image is delivered and the package declares no asset dependency.  Editorial formatting only, in addition: an AMS \texttt{amsart} preamble that restates the article's own declarations and macros used by the retained lines, namely the environments \texttt{theorem}, \texttt{lemma} on separate counters, together with the standard packages those lines need.  The retained environments print in this document as Theorem 1, Lemma 1 in order of appearance, whereas the article prints the same items with the numbering of its own full document.  The complete native source of the article as distributed in the arXiv e-print for this exact version is bundled unchanged as \texttt{sources/200358/source0.tex}.
\begin{lemma}[Quadratic layer second derivative]\label{lem:app-second-dist}
Let $L,D,W$ be $C^1$ functions on an interval
$I\subset(0,\infty)$, with $L>0$, $D\ge0$, and $W\ge0$. Assume that
the measures and integrals below are locally finite. Define, for $x>0$,
\[
 F(x)=\int_I W(R)\frac12\Big(
 (L(R)^2-(x+D(R))^2)_+
 +(L(R)^2-(x-D(R))^2)_+
 \Big)\,\dd R,
\]
and set $G(q)=F(\sqrt q)$ for $q>0$. Let
\[
 \nu_-=(|L-D|)_\#(W L\,\dd R),
 \qquad
 \nu_+=(L+D)_\#(W L\,\dd R)
\]
be the pushforward measures on $(0,\infty)$. Then, as a signed
distribution in the variable $x>0$,
\begin{equation}\label{eq:app-dist-measure}
 xF''-F'
 =
 x(\nu_-+\nu_+)
 -
 \left(
 \int_{\{|L-D|<x<L+D\}}W(R)D(R)\,\dd R
 \right)\dd x.
\end{equation}
Consequently, at every common regular value $x>0$ of $|L-D|$ and
$L+D$,
\begin{align}\label{eq:app-dist-lemma}
4x^3G''(x^2)=&\
 x\sum_{|L-D|=x}\frac{W(R)L(R)}{|(L-D)'(R)|}
 +x\sum_{L+D=x}\frac{W(R)L(R)}{|(L+D)'(R)|}\notag\\
&\quad
 -\int_{\{|L-D|<x<L+D\}}W(R)D(R)\,\dd R.
\end{align}
In particular, critical or colliding roots are already accounted for
by the pushforward measures in \eqref{eq:app-dist-measure}; no
separate choice of root branches is required.
\end{lemma}

\begin{equation}
 xF''-F'
 =
 x(\nu_-+\nu_+)
 -
 \left(
 \int_{\{|L-D|<x<L+D\}}W(R)D(R)\,\dd R
 \right)\dd x.
\end{equation}

\begin{theorem}[Dual layer theorem]\label{thm:app-dlt}
Let $1/2<\alpha<1$, let $\beta\ge1+2\alpha$, and let $A\ge0$. Set
\[
 h(s)=\frac{1+As}{(1-s^p)^\alpha},
 \qquad
 g(s)=\frac{1-As}{(1-s^p)^\alpha},
 \qquad 0\le s<1.
\]
Let $x>0$ be a common regular value of $|g|$ and $h$, and assume
that no corresponding root lies at an endpoint. Then
\begin{align}\label{eq:app-dlt}
 &x\sum_{|g|=x}\frac{(1-s^p)^{\beta+\alpha-1}}{|g'(s)|}
 +x\sum_{h=x}\frac{(1-s^p)^{\beta+\alpha-1}}{|h'(s)|} \\
 &\hspace{1in}\ge
 A\int_{\{|g|<x<h\}}s(1-s^p)^{\beta+\alpha-1}\,\dd s.
\end{align}
\end{theorem}

\begin{theorem}[Beta--Rademacher preservation]\label{thm:app-append}
Let $1/2\le\alpha<1$, let $\beta\ge1+2\alpha$, and let $c\ge0$. Let
$Y_1,Y_2$ be square-integrable real random variables. Let
$T\sim\operatorname{Beta}(\alpha,\beta)$ and let $\eps$ be a Rademacher
sign, with $(T,\eps)$ independent of $(Y_1,Y_2)$. If
\[
 Y_1^2\cx Y_2^2,
\]
then
\[
 \left((1-T)^\alpha Y_1+cT^\alpha\eps\right)^2
 \cx
 \left((1-T)^\alpha Y_2+cT^\alpha\eps\right)^2.
\]
\end{theorem}

\begin{proof}[Proof of Theorem~\ref{thm:app-append}]
We first assume $1/2<\alpha<1$. It suffices to check lower-hinge functions $\Phi_\rho(q)=(\rho^2-q)_+$ and affine functions. Affine functions are preserved because the transformation adds the same second-moment contribution to both sides. The case $\rho=0$ is trivial, since $\Phi_0(q)\equiv0$ on $[0,\infty)$; hence assume $\rho>0$.

For $q\ge0$ set
\[
 G(q)=\E_{T,\eps}\left(\rho^2-\left((1-T)^\alpha\sqrt q+cT^\alpha\eps\right)^2\right)_+.
\]
The Rademacher average makes this depend only on $q$. We prove that $G$ is convex. Put
\[
 R=\left(\frac{T}{1-T}\right)^\alpha.
\]
Then
\[
 (1-T)^\alpha=(1+R^p)^{-\alpha},
 \qquad
 T^\alpha=R(1+R^p)^{-\alpha},
\]
and the density of $R$ is a positive multiple of $(1+R^p)^{-\alpha-\beta}\dd R$. Hence, up to a positive constant,
\[
G(q)=\int_0^\infty (1+R^p)^{-\beta-3\alpha}\frac12\sum_{\pm}
\left(\rho^2(1+R^p)^{2\alpha}-(\sqrt q\pm cR)^2\right)_+\,\dd R.
\]
Set
\[
 L(R)=\rho(1+R^p)^\alpha,
 \qquad
 D(R)=cR,
 \qquad
 W(R)=(1+R^p)^{-\beta-3\alpha}.
\]
Lemma~\ref{lem:app-second-dist} gives a second derivative formula. Now pass to
\[
 s=\frac{R}{(1+R^p)^{1/p}},
 \qquad
 R=\frac{s}{(1-s^p)^\alpha},
 \qquad
 \dd R=(1-s^p)^{-\alpha-1}\dd s.
\]
Then
\[
 1+R^p=(1-s^p)^{-1},
 \qquad
 L(R)\pm D(R)=\rho\frac{1\pm(c/\rho)s}{(1-s^p)^\alpha}.
\]
Writing $A=c/\rho$ and $y=x/\rho$, the dual layer functions are
\[
 h(s)=\frac{1+As}{(1-s^p)^\alpha},
 \qquad
 g(s)=\frac{1-As}{(1-s^p)^\alpha}.
\]
The weights transform as follows. Since
\[
 W(R)L(R)=\rho(1+R^p)^{-\beta-2\alpha}
 =\rho(1-s^p)^{\beta+2\alpha}
\]
and
\[
 \frac{\dd R}{\dd s}=(1-s^p)^{-\alpha-1},
\]
we have, at endpoint roots,
\begin{align*}
 \frac{W(R)L(R)}{|(L+D)'(R)|}
 &=
 \frac{(1-s^p)^{\beta+\alpha-1}}{|h'(s)|},\\
 \frac{W(R)L(R)}{|(L-D)'(R)|}
 &=
 \frac{(1-s^p)^{\beta+\alpha-1}}{|g'(s)|}.
\end{align*}
For the bulk term,
\[
 W(R)D(R)\,\dd R
 =\rho A s(1-s^p)^{\beta+\alpha-1}\,\dd s.
\]
Thus, at every common regular value, the density of the signed
measure in \eqref{eq:app-dist-measure} is $\rho$ times the dual-layer
expression in Theorem~\ref{thm:app-dlt}, with parameter $A=c/\rho$
and threshold $y=x/\rho$. Indeed, $h'(s)>0$, while the sign of $g'(s)$ is the sign of
$s^{p-1}-A$. Thus $h$ is strictly increasing, and $g$ is strictly
monotone on each of at most two intervals. After splitting once more at
the possible zero of $g$, the same is true of $|g|$. On each monotone
branch, the one-dimensional change-of-variables formula shows that the
pushforward of the absolutely continuous endpoint measure is absolutely
continuous. The isolated critical values create no atoms. Consequently,
at almost every $x>0$ the Radon--Nikodym density is given by the coarea
sums above. The dual layer inequality shows that this density, and hence
the signed measure $xF''-F'$, is nonnegative. By the change of variables $q=x^2$, this is
equivalent to $G''\ge0$ distributionally on $(0,\infty)$. Continuity
at $0$ extends convexity to $[0,\infty)$.

Since $Y_1^2\cx Y_2^2$ and $G$ is convex,
\[
 \E G(Y_1^2)\le\E G(Y_2^2),
\]
which is the desired lower-hinge inequality after appending the coordinate.
This proves the theorem for $1/2<\alpha<1$.

It remains to treat $\alpha=1/2$. Choose $\alpha_j\downarrow1/2$ and
set
\[
 \beta_j=\max\{\beta,1+2\alpha_j\}.
\]
Then $\beta_j\to\beta$ and
$T_j\sim\operatorname{Beta}(\alpha_j,\beta_j)$ converges in law to
$T\sim\operatorname{Beta}(1/2,\beta)$. Apply the already proved case
to $(\alpha_j,\beta_j)$ and test with a lower hinge
$q\mapsto(\rho^2-q)_+$. Its value is bounded by $\rho^2$, so the
corresponding expectations converge as $j\to\infty$; the transformed
variables converge in distribution because the beta laws and the
exponents converge. The limiting lower-hinge inequality is the desired
one at $\alpha=1/2$. Equality of the transformed second moments follows
directly from $\E Y_1^2=\E Y_2^2$ and the symmetry of $\eps$. The
lower-hinge criterion completes the proof at the endpoint.
\end{proof}

\end{document}
